\VerdierPhysicalPageAnchor{0113}%
\begingroup
\small
\setlength{\abovedisplayskip}{4pt plus 2pt minus 2pt}%
\setlength{\belowdisplayskip}{4pt plus 2pt minus 2pt}%

\VerdierSemanticAnchor{ch2-1-1-5-additivity-announcement}
\noindent We shall show in (1.2.7) that cohomological functors are
additive and that multi-exact functors are multi-additive.

\VerdierSemanticAnchor{ch2-1-1-5-cohomology-notation}
Let $H\colon\mathcal D\longrightarrow\mathcal A$ be a cohomological
functor. We most often use the notation $H^0(X)=H(X)$,
$H^n(X)=H(X[n])$. Using axiom TRII, one sees that a cohomological functor
associates with every distinguished triangle

\VerdierSemanticAnchor{ch2-1-1-5-long-exact-setup}
\[
  X\xrightarrow{u}Y\xrightarrow{v}Z\xrightarrow{w}X[1]
\]
a long exact sequence:

\VerdierSemanticAnchor{ch2-1-1-5-1}
\begin{equation}
  \cdots\longrightarrow H^{-1}(Z)
  \xrightarrow{H^{-1}(-w)}H^0(X)
  \xrightarrow{H^0(u)}H^0(Y)
  \xrightarrow{H^0(v)}H^0(Z)
  \xrightarrow{H^0(w)}\cdots
  \tag{1.1.5.1}
\end{equation}

\VerdierSemanticAnchor{ch2-1-1-6}
\medskip
\noindent\textbf{1.1.6.}\quad
Let $\mathcal D$ be a triangulated category. We most often use the
following terminology: a triangle

\VerdierSemanticAnchor{ch2-1-1-6-antidistinguished-triangle}
\[
  X\xrightarrow{u}Y\xrightarrow{v}Z\xrightarrow{w}X[1]
\]
in $\mathcal D$ is called \emph{antidistinguished} if the triangle
\[
  X\xrightarrow{u}Y\xrightarrow{v}Z\xrightarrow{-w}X[1]
\]
is distinguished.

\VerdierSemanticAnchor{ch2-1-1-6-antiexact-functor}
A $\mathbb Z$-functor from a triangulated category $\mathcal D'$ to
$\mathcal D$ is called \emph{antiexact} if it carries the distinguished
triangles of $\mathcal D'$ to antidistinguished triangles of $\mathcal D$.

\VerdierSemanticAnchor{ch2-1-1-6-exact-functor-notation}
Most often, when denoting an exact functor from $\mathcal D'$ to
$\mathcal D$, we omit from the notation the isomorphisms commuting with
translations.

\VerdierSemanticAnchor{ch2-1-1-7}
\medskip
\noindent\textbf{1.1.7.}\quad
Let $\mathcal C$ be a strict $\mathbb Z$-category (Chapter I, 1.2.2),
with

\VerdierSemanticAnchor{ch2-1-1-7-translations}
\[
  X\longmapsto X[n],\qquad n\in\mathbb Z,
\]
the translation functors of $\mathcal C$. Let $\mathcal C^\circ$ be the
opposite $\mathbb Z$-category (Chapter I, 1.3.9).

\VerdierSemanticAnchor{ch2-1-1-7-opposite-morphism}
For every morphism $u$ of $\mathcal C$, denote by $((u))$ the
corresponding morphism of $\mathcal C^\circ$. This notation is abusive,
since the set of objects and the set of morphisms of $\mathcal C^\circ$
are equal respectively to the set of objects and the set of morphisms of
$\mathcal C$. Denote by

\VerdierSemanticAnchor{ch2-1-1-7-opposite-translations}
\[
  X\longmapsto X[[n]],\qquad n\in\mathbb Z,
\]
the translation functors of $\mathcal C^\circ$.

\VerdierSemanticAnchor{ch2-1-1-7-translation-identity}
Recall that, for every object $X$ of $\mathcal C$ and every integer
$n\in\mathbb Z$ (Chapter I, 1.3.9),
\[
  X[-n]=X[[n]].
\]

\VerdierSemanticAnchor{ch2-1-1-7-opposite-category}
The \emph{opposite triangulated category} of a triangulated category
$\mathcal D$ is defined to be the strict opposite $\mathbb Z$-category
$\mathcal D^\circ$ of the category $\mathcal D$ (Chapter I, 1.3.9),
equipped with the following class of distinguished triangles. A triangle
in $\mathcal D^\circ$

\VerdierSemanticAnchor{ch2-1-1-7-opposite-triangle}
\[
  X[[-1]]\xrightarrow{((w))}Z\xrightarrow{((v))}Y
  \xrightarrow{((u))}X
\]
is distinguished if and only if the corresponding triangle in the
category $\mathcal D$

\VerdierSemanticAnchor{ch2-1-1-7-corresponding-antitriangle}
\[
  X\xrightarrow{u}Y\xrightarrow{v}Z\xrightarrow{w}X[1]
\]
is \emph{antidistinguished}.

\VerdierSemanticAnchor{ch2-1-1-7-triangulated-axioms}
It is easily shown that the family of triangles of $\mathcal D^\circ$
defined in this way satisfies the axioms of triangulated categories.
\endgroup

